\documentclass[a4paper,12pt]{article}

\usepackage[T2A]{fontenc}
\usepackage[utf8]{inputenc}
\usepackage[russian]{babel}
\usepackage{amsmath,amssymb,mathtools}
\usepackage{helvet}
\renewcommand{\familydefault}{\sfdefault}
\usepackage{icomma}
\usepackage[a4paper,margin=2.15cm,headheight=16pt]{geometry}
\usepackage{microtype}
\usepackage{tikz}
\usetikzlibrary{calc,arrows.meta}
\usepackage{tabularx,booktabs}
\usepackage{enumitem}
\usepackage{hyperref}
\usepackage{parskip}
\usepackage{fancyhdr}
\usepackage{setspace}
\usepackage{xcolor}

% -------------------------------------------------
% Цветовая схема
% -------------------------------------------------
\definecolor{accent}{HTML}{008080}
\definecolor{darkaccent}{HTML}{004D4D}
\definecolor{textgray}{HTML}{333333}
\definecolor{softgray}{HTML}{EEF3F3}
\definecolor{linegray}{HTML}{B9C8C8}
\definecolor{warning}{HTML}{8B4B4B}

% -------------------------------------------------
% Общие настройки
% -------------------------------------------------
\onehalfspacing
\setlength{\parindent}{0pt}
\setlength{\emergencystretch}{2em}

\hypersetup{
    colorlinks=true,
    linkcolor=darkaccent,
    urlcolor=darkaccent
}

\setlist[itemize]{
    leftmargin=1.6em,
    itemsep=0.3em,
    topsep=0.3em
}

\setlist[enumerate]{
    leftmargin=2em,
    itemsep=0.7em,
    topsep=0.5em
}

\newcolumntype{Y}{>{\raggedright\arraybackslash}X}

\newcommand{\key}[1]{\textcolor{darkaccent}{\textbf{#1}}}
\newcommand{\warningtext}[1]{\textcolor{warning}{\textbf{#1}}}

% -------------------------------------------------
% Колонтитулы
% -------------------------------------------------
\pagestyle{fancy}
\fancyhf{}
\fancyhead[L]{\small Ксения Клыкова}
\fancyhead[R]{\small ЕГЭ: текстовые задачи}
\fancyfoot[C]{\small\thepage}
\renewcommand{\headrulewidth}{0.4pt}

% -------------------------------------------------
% Заголовки
% -------------------------------------------------
\usepackage{titlesec}

\titleformat{\section}
  {\Large\bfseries\color{darkaccent}}
  {}
  {0pt}
  {}
  [\vspace{0.15em}\color{linegray}\titlerule]

\titleformat{\subsection}
  {\large\bfseries\color{accent}}
  {}
  {0pt}
  {}

\titlespacing*{\section}{0pt}{1.4em}{0.8em}
\titlespacing*{\subsection}{0pt}{1.1em}{0.45em}

\begin{document}

% =================================================
% ТИТУЛЬНЫЙ ЛИСТ
% =================================================

\begin{titlepage}
\thispagestyle{empty}

\begin{center}

\vspace*{3.2cm}

{\Large\textcolor{accent}{ЕГЭ по математике}\par}

\vspace{1.2cm}

{\fontsize{30}{36}\selectfont
\bfseries
\textcolor{darkaccent}{Чек-лист}\par}

\vspace{0.7cm}

{\LARGE
\bfseries
Повторение текстовых задач\par}

\vspace{0.35cm}

{\large
работа, движение, протяжённые объекты,\\
формульные модели, смеси и сплавы\par}

\vspace{2.5cm}

{\large
Лёвин Артём Александрович\\[0.3em]
эксклюзивно для Ксении Клыковой\par}

\vfill

{\large \today\par}

\end{center}

\begin{tikzpicture}[remember picture,overlay]
\draw[
    accent!55,
    line width=1.3pt
]
([xshift=1.4cm,yshift=1.4cm]current page.south west)
.. controls
([xshift=5.0cm,yshift=2.0cm]current page.south west)
and
([xshift=-5.0cm,yshift=0.8cm]current page.south east)
..
([xshift=-1.4cm,yshift=1.4cm]current page.south east);
\end{tikzpicture}

\end{titlepage}

% =================================================
% ВВЕДЕНИЕ
% =================================================

\section*{Главная идея урока}

Большинство текстовых задач удобно начинать не с уравнения, а с
\key{правильного распределения данных по величинам}.

Для задач на движение обычно используются
\[
S=vt,
\qquad
t=\frac{S}{v},
\qquad
v=\frac{S}{t}.
\]

Для задач на работу аналогичная структура имеет вид
\[
A=pt,
\]
где \(A\) --- объём выполненной работы,
\(p\) --- производительность,
\(t\) --- время.

Для смесей и сплавов используется равенство
\[
m_{\text{вещества}}
=
m_{\text{смеси}}\cdot c,
\]
где \(c\) --- концентрация, записанная десятичной дробью.

\key{Перед составлением уравнения всегда определите:}
\begin{itemize}
    \item что означает каждая величина;
    \item в каких единицах она измеряется;
    \item что именно удобно обозначить неизвестной;
    \item какое физическое или математическое равенство связывает строки таблицы.
\end{itemize}

% =================================================
% 1. ЗАДАЧИ НА РАБОТУ
% =================================================

\section{Задачи на совместную работу}

\subsection*{Базовая модель}

Если весь заказ принимается за \(1\), то производительность исполнителя,
выполняющего заказ за \(t\) часов, равна
\[
p=\frac{1}{t}.
\]

Если два исполнителя работают одновременно, их производительности
\key{складываются}:
\[
p_{\text{общ}}=p_1+p_2.
\]

Поэтому
\[
\frac1{t_1}+\frac1{t_2}
=
\frac1{t_{\text{совм}}}.
\]

\begin{table}[ht]
\centering
\caption{Стандартная таблица для задачи на работу}
\renewcommand{\arraystretch}{1.35}
\begin{tabularx}{\linewidth}{Y Y Y Y}
\toprule
Исполнитель &
Производительность &
Время &
Работа \\
\midrule
Первый &
\(\dfrac1{t_1}\) &
\(t_1\) &
\(1\) \\
Второй &
\(\dfrac1{t_2}\) &
\(t_2\) &
\(1\) \\
Вместе &
\(\dfrac1{t_{\text{совм}}}\) &
\(t_{\text{совм}}\) &
\(1\) \\
\bottomrule
\end{tabularx}
\end{table}

\subsection*{Типовой пример}

Две машины, работая вместе, выполняют заказ за \(4\) часа.
Первая машина выполняет тот же заказ на \(6\) часов быстрее второй.

Обозначим через \(x\) время работы второй машины.

Тогда
\[
t_2=x,
\qquad
t_1=x-6.
\]

Здесь число \(6\) относится именно к \key{разности времени}, поэтому оно
должно появиться в столбце времени, а не производительности.

Получаем
\[
\frac1{x-6}+\frac1x=\frac14.
\]

\key{Ограничение:}
\[
x>6.
\]

Решаем:
\[
4x+4(x-6)=x(x-6),
\]
\[
x^2-14x+24=0,
\]
\[
x=12
\quad\text{или}\quad
x=2.
\]

Значение \(x=2\) не удовлетворяет условию \(x>6\), поэтому
\[
t_2=12\text{ ч},
\qquad
t_1=6\text{ ч}.
\]

\subsection*{Проверка}

\[
\frac16+\frac1{12}
=
\frac2{12}+\frac1{12}
=
\frac14.
\]

\key{Вывод:} если в условии сказано
«на столько-то часов быстрее», речь идёт о \textbf{времени},
а не о производительности.

% =================================================
% 2. ПРОТЯЖЕННЫЕ ОБЪЕКТЫ
% =================================================

\section{Движение протяжённых объектов}

В обычных задачах тело можно считать точкой.
Поезд, колонна или другой протяжённый объект имеет собственную длину,
которую необходимо учитывать.

\subsection*{Два поезда полностью проходят друг мимо друга}

Для полного прохождения необходимо преодолеть суммарную длину
\[
S=L_1+L_2.
\]

Если поезда движутся навстречу:
\[
v_{\text{отн}}=v_1+v_2.
\]

Если они движутся в одном направлении:
\[
v_{\text{отн}}=|v_1-v_2|.
\]

Следовательно,
\[
\boxed{
t=\frac{L_1+L_2}{v_{\text{отн}}}
}
\]

\begin{center}
\begin{tikzpicture}[>=Latex,scale=0.95]
    \draw[line width=0.7pt] (0,0)--(11,0);

    \draw[fill=softgray,draw=darkaccent,line width=0.8pt]
        (1,0.15) rectangle (3.3,0.75);
    \draw[fill=softgray,draw=darkaccent,line width=0.8pt]
        (7.7,0.15) rectangle (10.2,0.75);

    \node at (2.15,0.45) {\small поезд 1};
    \node at (8.95,0.45) {\small поезд 2};

    \draw[->,accent,line width=1.2pt] (3.5,0.45)--(5.2,0.45);
    \draw[->,accent,line width=1.2pt] (7.5,0.45)--(5.8,0.45);

    \node[below] at (2.15,0) {\small \(L_1\)};
    \node[below] at (8.95,0) {\small \(L_2\)};
\end{tikzpicture}
\end{center}

\subsection*{Пример: встречное движение}

Длины поездов:
\[
150\text{ м}
\quad\text{и}\quad
210\text{ м}.
\]

Скорости:
\[
72\text{ км/ч}
\quad\text{и}\quad
54\text{ км/ч}.
\]

Сначала переводим длины в километры:
\[
150\text{ м}=0{,}15\text{ км},
\qquad
210\text{ м}=0{,}21\text{ км}.
\]

Тогда
\[
t=
\frac{0{,}15+0{,}21}{72+54}
=
\frac{0{,}36}{126}
=
\frac1{350}\text{ ч}.
\]

Чтобы перейти из часов в секунды,
нужно умножить на \(3600\):
\[
t=
\frac{3600}{350}
=
\frac{72}{7}
\approx 10{,}29\text{ с}.
\]

\subsection*{Поезд и мост}

Если поезд длиной \(L_{\text{п}}\) полностью проходит мост длиной
\(L_{\text{м}}\), то передняя часть поезда должна пройти
\[
L_{\text{п}}+L_{\text{м}}.
\]

Следовательно,
\[
\boxed{
t=
\frac{L_{\text{п}}+L_{\text{м}}}{v}
}
\]

Мост неподвижен, поэтому его скорость равна нулю.

% =================================================
% 3. ЕДИНИЦЫ ИЗМЕРЕНИЯ
% =================================================

\section{Единицы измерения: обязательная проверка}

В формуле
\[
t=\frac{S}{v}
\]
единицы расстояния должны соответствовать единицам скорости.

Если скорость дана в \(\text{км/ч}\), удобно переводить метры в километры:
\[
1\text{ км}=1000\text{ м}.
\]

Для времени:
\[
1\text{ ч}=60\text{ мин}=3600\text{ с}.
\]

Поэтому
\[
t_{\text{мин}}=60t_{\text{ч}},
\qquad
t_{\text{с}}=3600t_{\text{ч}}.
\]

\warningtext{Типичная ошибка:}
домножить числитель и знаменатель дроби на \(60\).

Например,
\[
\frac{0{,}6}{90}
=
\frac{36}{5400}.
\]

Это та же самая дробь:
её числовое значение и единица измерения не изменились.

Чтобы перевести найденное время из часов в секунды, необходимо писать
\[
\frac{0{,}6}{90}\cdot3600.
\]

\key{Правило:}
перевод единиц изменяет величину числового значения,
а одинаковое умножение числителя и знаменателя только изменяет вид дроби.

% =================================================
% 4. ФОРМУЛЬНАЯ МОДЕЛЬ
% =================================================

\section{Задачи с заданной формулой}

Рассмотрим модель
\[
H=5t^2,
\]
где \(H\) --- расстояние до поверхности воды в метрах,
а \(t\) --- время падения камня в секундах.

До дождя
\[
t_1=1{,}5\text{ с}.
\]

После дождя уровень воды поднялся, поэтому расстояние до воды
\key{уменьшилось}. Следовательно, время падения также уменьшилось.

Если оно изменилось на \(0{,}1\) с, то
\[
t_2=1{,}4\text{ с}.
\]

Искомый подъём уровня воды равен
\[
\Delta H
=
H_{\text{до}}-H_{\text{после}}.
\]

Получаем
\[
\Delta H
=
5\cdot1{,}5^2-5\cdot1{,}4^2.
\]

Можно считать квадраты непосредственно, но быстрее применить
формулу разности квадратов:
\[
a^2-b^2=(a-b)(a+b).
\]

Тогда
\[
\Delta H
=
5(1{,}5^2-1{,}4^2)
=
5(1{,}5-1{,}4)(1{,}5+1{,}4),
\]
\[
\Delta H
=
5\cdot0{,}1\cdot2{,}9
=
1{,}45\text{ м}.
\]

\key{Полезная идея:}
если в формуле присутствуют квадраты близких чисел, сначала проверьте,
нельзя ли использовать разность квадратов.

% =================================================
% 5. СМЕСИ И СПЛАВЫ
% =================================================

\section{Смеси, растворы и сплавы}

Задачи на растворы и сплавы имеют одну математическую модель.

Главное равенство:
\[
m_{\text{вещества}}
=
m_{\text{смеси}}\cdot c.
\]

Процент обязательно переводим в десятичную дробь:
\[
10\%=0{,}10,
\qquad
35\%=0{,}35,
\qquad
30\%=0{,}30.
\]

\begin{table}[ht]
\centering
\caption{Стандартная таблица для смеси}
\renewcommand{\arraystretch}{1.4}
\begin{tabularx}{\linewidth}{Y Y Y Y}
\toprule
Смесь &
Масса смеси &
Концентрация &
Масса вещества \\
\midrule
Первая &
\(x\) &
\(c_1\) &
\(c_1x\) \\
Вторая &
\(y\) &
\(c_2\) &
\(c_2y\) \\
Полученная &
\(x+y\) &
\(c\) &
\(c(x+y)\) \\
\bottomrule
\end{tabularx}
\end{table}

\subsection*{Типовой пример со сплавами}

Первый сплав содержит \(10\%\) никеля,
второй --- \(35\%\).

Получили \(225\) кг сплава с содержанием никеля \(30\%\).

Пусть
\[
x
\]
--- масса первого сплава,
а
\[
y
\]
--- масса второго.

Из общей массы:
\[
x+y=225.
\]

Из массы никеля:
\[
0{,}10x+0{,}35y=0{,}30\cdot225.
\]

Получаем систему
\[
\begin{cases}
x+y=225,\\
0{,}10x+0{,}35y=67{,}5.
\end{cases}
\]

Её решение:
\[
x=45,
\qquad
y=180.
\]

Масса первого сплава меньше массы второго на
\[
180-45=135\text{ кг}.
\]

\key{Важно:}
если в задаче спрашивается разность масс, вовсе не обязательно
обозначать эту разность через \(x\).
Часто удобнее сначала найти сами массы.

% =================================================
% 6. ДВИЖЕНИЕ С РАЗВОРОТОМ
% =================================================

\section{Движение до точки и обратно}

Два человека одновременно вышли из дома к опушке леса,
расстояние до которой равно
\[
4{,}2\text{ км}.
\]

Скорости:
\[
v_{\text{быстр}}=4{,}5\text{ км/ч},
\qquad
v_{\text{медл}}=2{,}5\text{ км/ч}.
\]

Быстрый дошёл до опушки, развернулся и встретил медленного.

\begin{center}
\begin{tikzpicture}[>=Latex,scale=1]
    \draw[line width=0.8pt] (0,0)--(10,0);

    \fill[darkaccent] (0,0) circle (2pt);
    \fill[accent] (7.15,0) circle (2pt);
    \fill[darkaccent] (10,0) circle (2pt);

    \node[below] at (0,-0.05) {\small дом};
    \node[below] at (7.15,-0.05) {\small встреча};
    \node[below] at (10,-0.05) {\small опушка};

    \draw[<->,accent,line width=1pt]
        (7.15,0.75)--(10,0.75);
    \node[above] at (8.58,0.75) {\small \(x\)};

    \draw[->,darkaccent,line width=1pt]
        (0,1.45)--(9.7,1.45);
    \draw[->,darkaccent,line width=1pt]
        (9.7,1.45)--(7.3,1.45);

    \node[above] at (4.9,1.45)
        {\small путь быстрого};
\end{tikzpicture}
\end{center}

Обозначим через \(x\) расстояние от места встречи до опушки.

Тогда путь быстрого:
\[
S_1=4{,}2+x,
\]
а путь медленного:
\[
S_2=4{,}2-x.
\]

Так как они вышли одновременно и встретились в один момент,
времена движения равны:
\[
\frac{4{,}2+x}{4{,}5}
=
\frac{4{,}2-x}{2{,}5}.
\]

Получаем
\[
2{,}5(4{,}2+x)
=
4{,}5(4{,}2-x),
\]
\[
10{,}5+2{,}5x
=
18{,}9-4{,}5x,
\]
\[
7x=8{,}4,
\]
\[
x=1{,}2.
\]

Но \(x\) --- это не ответ задачи.

Расстояние от дома до места встречи:
\[
4{,}2-1{,}2=3\text{ км}.
\]

\key{Методический вывод:}
иногда удобная неизвестная не совпадает с тем, что требуется найти.
После решения уравнения обязательно ещё раз прочитайте вопрос задачи.

% =================================================
% 7. АЛЬТЕРНАТИВНЫЙ МЕТОД
% =================================================

\section{Одна неизвестная или система?}

Предыдущую задачу можно решить и с двумя неизвестными.

Пусть
\[
S_1
\]
--- путь быстрого,
а
\[
S_2
\]
--- путь медленного.

По равенству времени:
\[
\frac{S_1}{4{,}5}
=
\frac{S_2}{2{,}5}.
\]

Кроме того, суммарно два пути составляют удвоенное расстояние
от дома до опушки:
\[
S_1+S_2=2\cdot4{,}2=8{,}4.
\]

Получаем систему
\[
\begin{cases}
\dfrac{S_1}{4{,}5}=\dfrac{S_2}{2{,}5},\\[0.7em]
S_1+S_2=8{,}4.
\end{cases}
\]

Из первого уравнения:
\[
2{,}5S_1=4{,}5S_2,
\]
\[
5S_1=9S_2,
\]
\[
S_1=\frac95S_2.
\]

Подставляем:
\[
\frac95S_2+S_2=8{,}4,
\]
\[
14S_2=42,
\]
\[
S_2=3.
\]

Именно \(S_2\) является расстоянием от дома до места встречи.

\subsection*{Какой способ лучше?}

Оба решения корректны.

Однако на экзамене обычно выгоднее модель с
\key{одной неизвестной}, если она получается естественно:
\begin{itemize}
    \item меньше вычислений;
    \item ниже риск арифметической ошибки;
    \item быстрее получается уравнение;
    \item легче контролировать смысл величин.
\end{itemize}

Система полезна как альтернативный способ и как проверка понимания модели.

% =================================================
% 8. РАЗНОСТЬ ПУТЕЙ
% =================================================

\section{Когда неизвестные сокращаются}

Иногда в таблице возникают две переменные, но это ещё не означает,
что обязательно понадобится система.

Пусть два объекта одновременно движутся в одном направлении.

Время движения обозначим через \(x\).

Скорость медленного:
\[
y,
\]
скорость быстрого:
\[
y+a.
\]

Тогда пройденные пути равны
\[
yx
\quad\text{и}\quad
(y+a)x.
\]

Если расстояние между объектами стало равно \(d\), то
\[
(y+a)x-yx=d.
\]

После раскрытия скобок:
\[
yx+ax-yx=d,
\]
поэтому
\[
ax=d
\]
и
\[
\boxed{x=\frac da}.
\]

Переменная \(y\) сократилась.

\key{Вывод:}
две буквы в таблице не всегда означают две необходимые неизвестные
в итоговом уравнении.

% =================================================
% 9. АЛГОРИТМ
% =================================================

\section{Универсальный алгоритм решения}

\begin{enumerate}
    \item Внимательно прочитайте вопрос задачи.
    \item Определите тип модели:
    работа, движение, смесь, формула.
    \item Выпишите известные величины и их единицы.
    \item При необходимости сделайте схему или таблицу.
    \item Выберите неизвестную так, чтобы число переменных было минимальным.
    \item Выразите остальные величины через выбранную неизвестную.
    \item Найдите равенство, связывающее строки таблицы.
    \item Составьте уравнение.
    \item Укажите естественные ограничения:
    время \(>0\), скорость \(>0\), масса \(>0\), знаменатель \(\ne0\).
    \item Решите уравнение.
    \item Отбросьте значения, противоречащие смыслу задачи.
    \item Вернитесь к вопросу: найденная переменная может ещё не быть ответом.
    \item Проверьте единицы измерения.
\end{enumerate}

% =================================================
% 10. ТИПИЧНЫЕ ОШИБКИ
% =================================================

\section{Типичные ошибки}

\begin{itemize}
    \item \warningtext{Часы записаны не в тот столбец.}
    Если сказано «на \(6\) часов быстрее», это разность времени.

    \item \warningtext{При встречном движении вычитают скорости.}
    При сближении:
    \[
    v_{\text{отн}}=v_1+v_2.
    \]

    \item \warningtext{При движении в одном направлении складывают скорости.}
    Для обгона:
    \[
    v_{\text{отн}}=|v_1-v_2|.
    \]

    \item \warningtext{Не учитывают длину поезда.}
    Для полного прохождения требуется пройти сумму соответствующих длин.

    \item \warningtext{Метры смешивают с километрами.}
    Перед подстановкой единицы необходимо согласовать.

    \item \warningtext{Часы переводят в секунды умножением на \(60\).}
    Верно:
    \[
    1\text{ ч}=3600\text{ с}.
    \]

    \item \warningtext{Домножение числителя и знаменателя принимают за перевод единиц.}
    Такое преобразование не изменяет значение дроби.

    \item \warningtext{После дождя выбирают большее время падения камня.}
    При поднявшемся уровне воды расстояние до воды уменьшается.

    \item \warningtext{Проценты не переводят в десятичные дроби.}
    Например:
    \[
    35\%=0{,}35.
    \]

    \item \warningtext{Получили \(x\) и сразу записали его в ответ.}
    Сначала проверьте, что именно обозначает \(x\).

    \item \warningtext{Слишком много неизвестных.}
    Если можно построить корректную модель с одной переменной,
    обычно она предпочтительнее.
\end{itemize}

% =================================================
% КРАТКАЯ ПАМЯТКА
% =================================================

\section{Краткая памятка перед экзаменом}

\begin{table}[ht]
\centering
\renewcommand{\arraystretch}{1.45}
\begin{tabularx}{\linewidth}{Y Y}
\toprule
Ситуация & Формула или ключевая связь \\
\midrule
Движение &
\(\displaystyle S=vt\) \\

Работа &
\(\displaystyle A=pt\) \\

Весь заказ выполнен за \(t\) &
\(\displaystyle p=\frac1t\) \\

Совместная работа &
\(\displaystyle p=p_1+p_2\) \\

Поезда навстречу &
\(\displaystyle
t=\frac{L_1+L_2}{v_1+v_2}\) \\

Поезда в одном направлении &
\(\displaystyle
t=\frac{L_1+L_2}{|v_1-v_2|}\) \\

Поезд проходит мост &
\(\displaystyle
t=\frac{L_{\text{поезда}}+L_{\text{моста}}}{v}\) \\

Смеси и сплавы &
\(\displaystyle m_{\text{вещества}}=cm_{\text{смеси}}\) \\

Разность квадратов &
\(\displaystyle a^2-b^2=(a-b)(a+b)\) \\

Часы \(\to\) минуты &
умножить на \(60\) \\

Часы \(\to\) секунды &
умножить на \(3600\) \\
\bottomrule
\end{tabularx}
\end{table}

% =================================================
% ТРЕНИРОВОЧНЫЙ БЛОК
% =================================================

\clearpage

\section*{Тренировочный блок}

Решите задачи самостоятельно.
В тестовой части экзамена важен прежде всего корректный ответ,
но вычисления должны оставаться математически последовательными.

\subsection*{Средний уровень}

\begin{enumerate}[label=\textbf{\arabic*.}]

\item
Две машины, работая вместе, выполняют заказ за \(6\) часов.
Первая машина выполняет этот заказ на \(5\) часов быстрее второй.
За сколько часов каждая машина выполнит заказ отдельно?

\item
Два поезда длиной \(180\) м и \(120\) м движутся навстречу друг другу
со скоростями \(54\) км/ч и \(36\) км/ч.
За сколько секунд они полностью пройдут друг мимо друга?

\item
Поезд длиной \(180\) м движется со скоростью \(72\) км/ч.
За сколько секунд он полностью пройдёт мост длиной \(420\) м?

\end{enumerate}

\subsection*{Выше среднего}

\begin{enumerate}[label=\textbf{\arabic*.},start=4]

\item
Два поезда длиной \(240\) м и \(160\) м движутся в одном направлении
со скоростями \(90\) км/ч и \(72\) км/ч.
Быстрый поезд догоняет медленный.
За сколько секунд он полностью пройдёт мимо него?

\item
Расстояние до поверхности воды в колодце вычисляется по формуле
\[
H=5t^2,
\]
где \(H\) измеряется в метрах, а \(t\) --- в секундах.
До дождя время падения камня составляло \(1{,}7\) с,
а после дождя уменьшилось на \(0{,}2\) с.
На сколько метров поднялся уровень воды?

\item
Смешали два сплава.
Первый содержит \(20\%\) никеля, второй --- \(50\%\).
Получили \(30\) кг сплава, содержащего \(38\%\) никеля.
На сколько килограммов масса первого сплава была меньше массы второго?

\item
От дома до опушки леса \(5{,}6\) км.
Два туриста одновременно вышли из дома в сторону опушки
со скоростями \(6\) км/ч и \(4\) км/ч.
Быстрый турист дошёл до опушки, сразу повернул обратно
и встретил медленного.
На каком расстоянии от дома произошла встреча?

\item
Два велосипедиста одновременно выехали из одной точки
в одном направлении.
Скорость первого на \(4\) км/ч больше скорости второго.
Через какое время расстояние между ними станет равным \(600\) м?
Ответ дайте в минутах.

\item
Первая машина выполняет заказ за \(9\) часов.
Работая вместе со второй машиной, они выполняют этот заказ за \(6\) часов.
За сколько часов этот заказ выполняет вторая машина отдельно?

\end{enumerate}

\subsection*{Сложная задача}

\begin{enumerate}[label=\textbf{\arabic*.},start=10]

\item
От дома до опушки леса \(7{,}2\) км.
Два туриста одновременно вышли из дома в сторону опушки.
Скорость быстрого туриста на \(2\) км/ч больше скорости медленного.

Быстрый турист дошёл до опушки, отдыхал там \(6\) минут,
после чего пошёл обратно.
Туристы встретились в \(1{,}2\) км от опушки.

Найдите скорости обоих туристов.

\end{enumerate}

% =================================================
% ОТВЕТЫ
% =================================================

\clearpage

\section*{Ответы к тренировочному блоку}

\begin{table}[ht]
\centering
\caption{Ответы}
\renewcommand{\arraystretch}{1.55}
\begin{tabularx}{\linewidth}{c Y}
\toprule
№ & Ответ \\
\midrule

1 &
Первая машина: \(10\) ч;
вторая машина: \(15\) ч. \\

2 &
\(12\) с. \\

3 &
\(30\) с. \\

4 &
\(80\) с. \\

5 &
\(3{,}2\) м. \\

6 &
На \(6\) кг. \\

7 &
\(4{,}48\) км от дома. \\

8 &
\(9\) мин. \\

9 &
\(18\) ч. \\

10 &
Медленный турист: \(4\) км/ч;
быстрый турист: \(6\) км/ч. \\

\bottomrule
\end{tabularx}
\end{table}

\vfill

\begin{center}
\textcolor{darkaccent}{
\textbf{Перед записью ответа всегда проверьте смысл найденной величины
и единицы измерения.}
}
\end{center}

\end{document}